\subsection{Riesgos del desarrollo del proyecto}\label{sec:sd8-desarrollo}
La perspectiva de desarrollo de T-22 evalúa la construcción y recepción de resultados, mientras la de implantación evalúa el cambio efectivo de operación. Ambas emplean las fichas y valoraciones vigentes de T-16: un escenario no crea otra puntuación ni otra reserva. El registro \path{componentes/escenarios_desarrollo_implantacion.csv} enlaza cada escenario con causa, consecuencia, evaluación y respuesta de esas fichas. Las actividades y pruebas descritas son compromisos futuros. \parencites[Formulario T-22; arts.~17--18; pp.~12--14, 68]{sd8-ba}[8.2--8.3]{sd8-ex}

\paragraph{Alcance y conocimiento sectorial: R-03/R-04/R-05/R-11.}
Una regla de ocupación, regreso, autorización escolar, reparto de consumo o restricción de acceso no acordada puede producir código consistente con un supuesto incorrecto y rechazo posterior. Análisis Funcional conserva decisiones DEC de SD3, supuesto, autoridad competente, casos de aceptación y dependencias antes de construir; Arquitectura verifica la selección física de SD4. La divergencia entre tabletas de SD3 y terminal seleccionado en SD4 sigue abierta. La falta de acuerdo bloquea la función dependiente y se eleva a Proyecto, sin completar datos ni imponer una decisión societaria mediante software. Los cambios se tramitan por CC-6 de T-9, con impacto y aprobación previa, conforme a BA 72 y RT-19.03; tampoco se cambia unilateralmente proveedor, región o mecanismo contractual. El límite acumulado del 20\,\% del valor original y la prohibición de cambiar objeto, naturaleza, garantías mínimas o calendario obligatorio siguen vigentes.

La brecha sectorial puede invalidar requisitos de náutica deportiva, zarpe, combustible, instalaciones flotantes, residuos o cobro de servicios. Proyecto coordinará revisión por especialistas competentes con Análisis Funcional y Arquitectura antes de aprobar reglas o especificaciones. Se enlaza a 1.3 para reglas y normativa, 1.7 para instalación, 1.8 para interpretación histórica y 1.9 para indicadores/pilotos; los referentes funcionales aportan operación y no sustituyen competencia normativa. Cada revisión producirá fuente oficial vigente, regla aplicable, decisión, caso de prueba y observaciones resueltas. Las identidades, disponibilidad y HH de estos especialistas deben incorporarse a T-15: su dedicación se estima por tarea y disponibilidad confirmada, sin presumir contratación. La planificación debe acreditar capacidad para ejecutar la acción. \parencite[numeral 16.2; RT-15.02; caps.~11--12]{sd8-c07}

\paragraph{Equipo, capacidad y productividad: R-04/R-05/R-13.}
Ausencia o rotación de un integrante de Synaptix puede detener integración, seguridad, pruebas o la coordinación conjunta Implantación/SRE. Proyecto distribuirá documentación, decisiones, repositorios y runbooks; verificará que otra persona competente pueda ejecutar el procedimiento, con permisos y agenda. El reemplazo del equipo clave requiere autorización previa y escrita del CLIENTE, perfil igual o superior acreditado, traslape mínimo de quince días hábiles y ningún costo adicional. La rotación no autorizada o superior al 30\,\% del equipo clave en doce meses constituye incumplimiento; el relevo de emergencia no elimina estas obligaciones. Ante ausencia abrupta se protege el trabajo y se suspende la actividad sin cobertura habilitada, mientras se tramita la sustitución. \parencite[art.~76; p.~39]{sd8-ba}

La existencia de roles no prueba disponibilidad para construir E2 mientras E1 opera. T-15 registra 936 HH de ingeniería de migración y 1.644 HH de ensayos/implantación de datos, con fronteras distintas; tampoco cubren toda la ejecución del contrato. En meses 13--15 ese último componente suma 232/96/96 HH, con 80/48/48 HH conjuntas de Operación e Implantación; en 19--20 suma 164/96 HH y 64/48 HH conjuntas. Son cargas propuestas de los componentes de SD7, no capacidad libre ni turnos completos. Proyecto sumará construcción E2, integraciones, seguridad, pruebas, atención E1, revisiones y acompañamiento por persona/período, descontando ausencias y considerando una sola agenda para José Lara. Las 328 HH iniciales de control de riesgos tampoco sustituyen esa nivelación. La discrepancia entre los antecedentes funcionales y la selección de 12.708 HH de SD7 sigue requiriendo conciliación; no se habilita un frente por disponer de un cargo en organigrama.

Contratos de integración, versiones o ambientes inconsistentes originan retrabajo, deuda, pérdida de productividad y defectos tardíos en 1.3--1.6/1.10. Integración y Desarrollo aplicarán el flujo de SD6/T-10: contrato versionado, cambios de esquema compatibles, integración incremental, automatización y evidencia asociada al artefacto. Calidad aplicará SD9 con casos obligatorios, regresión y controles de datos; un caso bloqueado o no ejecutado no cuenta como aprobado. Se limita trabajo en curso y se corrige en el paquete causante antes de promoción; una deuda que afecta seguridad, integridad o recepción no se difiere como mejora opcional. \parencite[19.2; cap.~20; pp.~33--35]{sd8-bt}

\paragraph{Verificación y recepción: R-09/R-10/R-12/R-04.}
Un defecto bloqueante o hallazgo crítico/alto invalida la puerta aunque el desarrollo esté terminado. Seguridad coordina remediación y Calidad repite casos afectados y regresión sobre la versión corregida, incluidos reensayos de intrusión cuando corresponda. La reserva debe incluir diagnóstico, corrección, revisión, despliegue en Preproducción y nueva verificación antes del hito; los seis días adicionales de PL-R02 y los plazos máximos de vulnerabilidades no demuestran por sí solos que esa cadena quepa. Proyecto debe insertar tareas, HH y recursos en T-14/T-15/T-18, conservando esta cabida pendiente. No se libera por agotamiento de reserva. \parencites[art.~17.3; p.~13]{sd8-ba}[cap.~20; pp.~34--35]{sd8-bt}

La carga de la marina, la ausencia de referentes o una decisión societaria sin resolver pueden impedir validación y firma. Proyecto acuerda desde arranque referente/suplente, agenda por proceso y expediente; registra consulta, recepción, observación y decisión. BA 18.3 reserva diez días hábiles al CLIENTE para revisar y diez días hábiles a Synaptix para subsanar cuando haya observaciones. Como sensibilidad, una primera revisión con observaciones añade hasta veinte días hábiles al tramo considerado, además de la verificación técnica propia; no es una probabilidad ni autorización de aceptación al día veinte. La segunda presentación con observaciones de la misma naturaleza se considera atraso imputable. Estos derechos no se reducen para ajustar H3, H5, H10 o el acta E2; su inserción conjunta en la red con recursos sigue pendiente. Sin aceptación expresa se bloquea el cambio oficial. \parencite[arts.~18 y 71--72; pp.~13--14, 37--38]{sd8-ba}

\subsection{Dependencias de adquisición, recinto e interfaces}\label{sec:sd8-dependencias-implantacion}
El CLIENTE compra hardware y ejecuta las obras excluidas; Synaptix especifica cantidades, variantes, interfaces, montaje y condiciones de recepción. Entrega tardía, cantidad equivocada o equipo no conforme impiden integración y pruebas. No se convierte la exclusión en ausencia de responsabilidad de coordinación. Proyecto mantiene las habilitaciones EXT de SD7 y Arquitectura verifica evidencia antes de aceptar el consumidor. La Tabla~\ref{tab:sd8-dependencias-p05} declara puertas necesarias, disparador, efecto y alternativa. Las jornadas son posiciones tempranas del cálculo CPM de SD7, no fechas contractuales confirmadas ni compromisos recibidos de terceros; su traducción a calendario y su confirmación siguen pendientes.

\begingroup\setlength{\tabcolsep}{3pt}
\begin{longtable}{L{\dimexpr.24\textwidth-2\tabcolsep\relax}L{\dimexpr.34\textwidth-2\tabcolsep\relax}L{\dimexpr.42\textwidth-2\tabcolsep\relax}}
\caption{Coordinación de dependencias críticas y alternativas permitidas}\label{tab:sd8-dependencias-p05}\\\hline
\EncabezadoTabla{Dependencia y coordinación} & \EncabezadoTabla{Puerta necesaria y disparador} & \EncabezadoTabla{Efecto y respuesta}\\\hline\endfirsthead
\hline\EncabezadoTabla{Dependencia y coordinación} & \EncabezadoTabla{Puerta necesaria y disparador} & \EncabezadoTabla{Efecto y respuesta}\\\hline\endhead
Equipos, red y telemetría; Proyecto y Arquitectura; R-06/11/01. & EXT-EQUIPOS, EXT-RED y EXT-TELEMETRIA: jornada 72 de la red; antes de 1.7 y pruebas consumidoras. Faltante, cantidad o variante errónea o recepción no conforme. & Afecta 1.2/1.7/1.10 y H3; recibir por lote conforme, preparar pruebas aisladas y rechazar lo inseguro. Simulador no acredita recepción física; sustitución sujeta a aprobación y nueva prueba.\\\hline
Obra civil, canalización, electricidad, torretas/postación; Proyecto y Arquitectura; R-01/11. & EXT-RECINTO: jornada 72; especificación y sitio seguro recibidos antes de montaje. Obra sin recepción, energía/protección/acceso insuficiente. & Afecta 1.2.3/1.7.2 y ventana marina. Preparar fuera del recinto sólo cuando sea admisible; reprogramar ventana autorizada y conservar prueba real de integración.\\\hline
Identidad y licencias; Proyecto e Integración; R-06/07/09. & EXT-IDENTIDAD y EXT-LICENCIAS: jornada 61; permisos y versiones habilitados antes del consumidor. Acceso o licencia no disponible. & Bloquea integración y pruebas negativas; ambiente aislado para construcción, sin habilitar producción ni cambiar región o proveedor para eludir la dependencia.\\\hline
Contabilidad, pagos y firma; Proyecto e Integración; R-03/12/04. & EXT-CONTABLE: jornada 274 del cálculo; contrato, credenciales y pruebas antes de 1.10/1.11 administrativos. No respuesta, aprobación pendiente o efecto incierto. & Afecta saldos y corte E2. Consulta por ID y conciliación; mantener autoridad vigente. El diario preserva hechos pero no inventa autorización de pago o firma.\\\hline
Autoridad, comunicaciones y gestor de residuos; Proyecto y referente funcional; R-03/11/15. & Fecha necesaria: antes del ensayo, piloto u ola que consume RT-05.23; fecha civil por confirmar. Canal, autorización o gestor no habilitado activa atraso. & Afecta 1.3/1.9/1.10/1.11 y evidencia ambiental. Canal oficial o registro asistido permitido con trazabilidad; no emitir zarpe, declaración ni retiro en nombre del tercero sin competencia.\\\hline
Permisos, acreditación y técnicos competentes; Proyecto y Arquitectura; R-05/11/13. & Antes de cada faena: acceso, formación, técnico y aptitud de variante/montaje marino o zona clasificada acreditados. Falta cualquiera de ellos. & Bloquea instalación y recepción; conservar montaje preparado en sitio seguro y obtener habilitación. Ni urgencia ni reserva autorizan equipo general en surtidor clasificado.\\\hline
\end{longtable}\endgroup
\FuenteTabla{SD7, componentes/cpm\_implementacion.csv (EXT, posición temprana); Caso 07, caps.~5--6 y 11, restricciones 9--10 y RT-05.23; BTT, cap.~8.}
El registro de dependencia conservará responsable externo designado, coordinador Synaptix, fecha civil requerida, evidencia, consumidor y última fecha admisible. Si la fecha confirmada excede la requerida o falta evidencia al iniciar el consumidor, Proyecto activa el riesgo y propaga el atraso a la EDT y la red; no consume una holgura de otra rama. Las respuestas técnicas de la sección~\ref{sec:sd8-tecnicos} y CC-6 aplican a cualquier alternativa que altere la oferta.

\subsection{Ventanas, cancelaciones y evidencia antes de temporada}\label{sec:sd8-ventanas-p05}
Se conserva desarrollo E1 1--12, marcha blanca 13--15 y producción 16; desarrollo E2 13--18, marcha blanca 19--20, producción 21 y Operación 21--56. El calendario civil de SD7 parte del supuesto 1 de diciembre de 2026, no de una fecha contractual confirmada. Las olas E1 propuestas para 1--14 de diciembre de 2027 deben caber completas antes del congelamiento; ningún mes relativo demuestra por sí solo compatibilidad. Proyecto validará el inicio y regenerará ventanas y reservas si cambia. \parencite[art.~17; pp.~12--13]{sd8-ba}

R-01/R-11/R-14 incorporan prohibición de intervenir sistemas y capacitar del 15 de diciembre al 15 de marzo y las catorce fechas náuticas; las campañas abril--mayo y septiembre--noviembre restringen lo que afecta al varadero, sin cerrar automáticamente todo el recinto. Junio--agosto facilita intervenciones en tierra pero no elimina temporales ni marejadas. Un aviso oficial con 24--48 horas suspende toda actividad programada, incluidos frentes remotos: no se ofrece seguir programando construcción durante el aviso como recuperación de plazo. Proyecto registra suspensión, capacidad afectada y condición autorizada de reanudación. \parencite[restricciones 11--12; numeral 13.2; Anexo B.2]{sd8-c07}

La reserva propuesta mantiene diez días hábiles por campaña marina, ya imputados en T-18 y contados una sola vez. Su base es un supuesto de planificación, no un histórico de cancelaciones. Como sensibilidad de una campaña se consideran cinco, diez y quince días hábiles suspendidos; el exceso sobre diez resulta cero, cero y cinco. Para H3 con margen 2,33 días, el último escenario arroja déficit 2,67 bajo la misma cadena; para H5 con margen once no genera déficit aritmético y deja seis, sujeto a ventanas y recursos. La sección~\ref{sec:sd8-contraste-ventanas} desarrolla la regla sin doble descuento. El número de episodios, duración y tiempo hasta recuperar acceso deben registrarse como supuestos; dos episodios de cinco días en una misma campaña consumen sus diez días, no generan veinte días de reserva. Varias campañas o suspensión de otro frente exigen recalcular la red conjunta y nivelar recursos. Esta comprobación sigue abierta: la respuesta exigida es absorber sin mover hitos, pero los escenarios adversos muestran que todavía no está acreditada para toda cancelación. \parencite[numeral 17.5; introducción y cap.~19]{sd8-c07}

Medidores se instalan sólo con embarcación desconectada o fuera del agua. Arquitectura coordina asociación amarra/medidor, torretas, protección y prueba antes de reconectar, en ventana permitida; no se habilita trabajo en varadero congelado aprovechando una campaña. R-15/R-IN5-01 consideran medición tardía o defectuosa que deja insuficiente serie antes de temporada 2029. E1 incluye medición inicial; E2 completa ambiente. Proyecto y referente ambiental deben fijar comienzo de serie válida, cobertura, continuidad y fecha de evaluación aceptable para el esquema competente; empezar en 2029 no satisface evidencia previa. La cabida con entrega de equipos y campañas no se declara comprobada. \parencite[numerales 13.1--13.2 y 17.3--17.5; Anexo B.2]{sd8-c07}

R-01/R-11/R-13 cubren instalación, mantenimiento o reposición aplazados en las cuarenta boyas a seis millas. Sin energía, datos ni personal, el acceso requiere embarcación y mar favorable. SRE coordina salida, personal competente, repuesto y retorno con operación marítima; ante falta de ventana mantiene estado e inspección pendientes explícitos y reprograma acceso seguro. El trayecto terrestre no sustituye acceso por mar ni una inspección prevista prueba disponibilidad continua. Plazo/costo logístico se incorporan al consumidor de SD7 antes de comprometer una intervención. \parencite[RT-21.16]{sd8-c07}

\subsection{Migración del legado y ausencia anterior al corte}\label{sec:sd8-legado-p05}
R-04/R-12 consideran ausencia de la única operadora del sistema de 2014 antes de extracción, cierre, interpretación de saldos, validación o corte administrativo E2. Prolongar convivencia no resuelve la ausencia de quien sostiene el origen. Desde el mes 6, Datos coordina inventario y copia recuperable, extracción documentada, diccionario de reportes y cierre mensual asistido; se conservan las dos sesiones iniciales de noventa minutos propuestas en SD5 y se distribuye el conocimiento. Administración designará un suplente funcional y Proyecto verificará su agenda antes del siguiente cierre, sin esperar al mes 21. Datos y su suplente técnico deben demostrar acceso y recuperación; Calidad verificará que el suplente funcional repita el procedimiento sin la operadora. La identidad, disponibilidad y demostración siguen pendientes: no se presume reemplazo existente.

Ante ausencia previa a esa demostración se preserva el origen y se activa registro controlado de hechos administrativos nuevos bajo responsable funcional designado, con ID, comprobante, fecha, autorización y conciliación diaria. Datos obtiene extracción de sólo lectura y soporte técnico habilitado, sin fabricar interpretación de saldos ni convertir ese apoyo en proveedor original. Si no hay persona autorizada para validar, el saldo y la decisión permanecen pendientes; Proyecto escala y bloquea migración/aceptación. Los procesos locales críticos continúan según BCP y autoridad de SD4/SD5. Esta contingencia requiere habilitación anticipada de personas, permisos y diario, y no depende de recuperar a la operadora para registrar hechos nuevos. \parencite[cap.~5; RT-05.15; numeral 17.6]{sd8-c07}

Corrupción, extracción parcial, doce años de origen sin documentación, documentos ilegibles o transformación incorrecta afectan 1.8.1--1.8.6/1.10.1/1.11.1. Datos conserva originales y huellas, perfila por conjunto, registra defecto y decisión, prueba reglas y archivo consultable independiente. Los 60 GB iniciales y 120 GB de sensibilidad de SD5 son hipótesis, no inventario ni límite. La ingeniería de herramientas (936 HH) no absorbe revisión documental completa, digitalización, ensayos o corte. La cobertura conserva padrón completo de socios y armadores, contratos vigentes completos con documentos, seis años de pagos/saldos, los catorce expedientes de embarcaciones abandonadas íntegros, registro completo de embarcaciones/invernadas y tres años de matrícula escolar. Historia no migrada se conserva consultable conforme a retención aplicable; reducir volumen no autoriza omitir conjuntos. \parencites[RT-05.11--15; p.~12]{sd8-bt}[RT-05.15]{sd8-c07}

Datos y Calidad efectuarán dos ensayos completos en Preproducción antes de cada migración definitiva, con versión de herramientas, manifiestos, tiempos medidos por fase y totales, conciliación y retorno probado. Se concilian recuentos, claves, sumas por moneda/período, referencias y documentos; expedientes y autorizaciones vigentes se revisan completos y los atributos no críticos usan el muestreo dirigido de SD5. Diferencia monetaria de carga, omisión no explicada, referencia crítica inválida, duplicado o pérdida detienen avance. Toda diferencia se explica y resuelve con validador funcional antes de aceptar; compensar diferencias de signo contrario no acredita exactitud. Durante convivencia se concilia diariamente y antes de cada decisión de corte; ningún ensayo se declara ejecutado en esta oferta.

El corte propuesto de SD5/SD7 es 22:00--02:00 en fecha habilitada; decisión antes de 00:30 y noventa minutos exclusivos para retorno previo a nuevas escrituras. Carga atrasada no consume esa reserva. Después de escrituras nuevas se conserva diario y respuestas externas, se cerca la autoridad fallida y se concilian/reaplican hechos sin repetir cargos, pagos o despachos. La estimación posterior de 120 minutos excede en treinta la reserva previa de noventa: son escenarios distintos y no se promete que ambos caben en la misma reserva. Si el origen no admite hechos nuevos o no se puede conciliar, el diario de contingencia queda como única autoridad administrativa hasta corrección controlada; no se restaura una imagen descartando hechos. Calidad mide el retorno completo y Proyecto exige ajuste de precarga/herramientas y ventana antes de habilitar el corte. \parencite[numeral 17.6]{sd8-c07}

\subsection{Riesgos de implantación y marcha blanca}\label{sec:sd8-implantacion-p05}
R-03/R-11/R-12/R-14 evalúan una puesta en servicio que falle aunque el software funcione aislado. Despliegue masivo simultáneo propaga un error a torre, zarpes, escuela, combustible y administración y puede impedir retorno. Implantación conserva las olas por proceso/zona de SD7, dependencias y autoridad de SD5: no cambia registro oficial de toda la marina en una sola activación. Cada ola requiere recepción del alcance previo, usuarios preparados, pruebas aceptadas, datos conciliados y retorno disponible. Los noventa/120 minutos corresponden al corte administrativo, no al retorno de toda la plataforma; SRE y Calidad deben medir recuperación de cada ola y función dentro de los objetivos BCP antes de avanzar. \parencites[RT-20.01--02; p.~35]{sd8-bt}[numeral 13.3]{sd8-c07}

El 22\,\% de amarras en tramo incorrecto es base histórica del riesgo de inventario, no una tasa atribuida a la carga nueva. Antes de cortar, referente de operaciones y Datos contrastan en terreno embarcación, eslora y demás dimensiones relevantes, puesto, estado y vínculo contractual, con discrepancia y decisión documentadas. Se conservan los permisos para cualquier movimiento físico. Un puesto ocupado no se libera por ausencia de señal y un expediente no se corrige por inferencia sin evidencia. Una diferencia crítica bloquea la ola y se concilia sin alterar arbitrariamente cobro o custodia. \parencite[numeral 17.6, punto 3]{sd8-c07}

Durante convivencia E1/E2, R-03/R-12 preservan una autoridad de escritura por dato y período, y contratos compatibles para los datos compartidos. Antes del corte administrativo el legado conserva su autoridad; la consulta contractual para E1 no transfiere esa escritura. Se comparan extracciones y registros controlados; no se exige doble digitación ordinaria. Una captura de contingencia excepcional conserva fuente/ID y se incorpora una sola vez con conciliación. E2 no degrada integridad, disponibilidad ni tiempos de E1. \parencites[art.~17.2; p.~13]{sd8-ba}[RT-20.03; p.~35]{sd8-bt}

R-04/R-08/R-09/R-12/R-14 consideran marcha blanca con pocas operaciones reales, mezcla insuficiente, indicador fuera de umbral, diferencia inexplicada, incidente abierto o persona sin capacitación/certificación exigida. Las cuatro últimas semanas deben cubrir volumen real comprometido e indicadores sostenidos; carga sintética no reemplaza esa operación. Calidad contrasta los mínimos y denominadores del plan SD7, disponibilidad/tiempos y conciliación diaria; Implantación verifica preparación individual. El cierre exige conjuntamente cero incidentes críticos/altos abiertos atribuibles, volumen real e indicadores durante esas cuatro semanas, ninguna diferencia no explicada, personal capacitado y certificado conforme al plan aprobado y acta de Contraparte. Si falla una condición, Synaptix extiende marcha blanca a su costo, sin recargo ni desplazamiento de fechas siguientes, quedando expuesta a multas de atraso. Proyecto no autoriza producción unilateralmente ni contabiliza semanas previas a una corrección como evidencia sostenida de la nueva versión sin comprobar su validez. \parencite[art.~17.3; p.~13]{sd8-ba}

En una mañana peak, una reversión que deje datos escolares o de regreso en una imagen anterior expone personas y custodia. SRE mantiene las funciones críticas locales de SD4/SD5 y el BCP de esta oferta; el retorno administrativo no revierte hechos operacionales ya confirmados. Ante pérdida de integridad se detiene la función insegura, se comunica por canal competente y se conserva captura segura, identificada y conciliable; no se promete continuidad electrónica cuando sólo existe contingencia manual. Los diarios separan zarpe, salida escolar, despacho y regreso con identidad, autorización y tiempo; el silencio no confirma regreso. Calidad prueba retorno con esa mezcla y nuevas operaciones, incluidos efectos externos inciertos, sin eludir competencias de autoridad o reglas de seguridad física.

R-05/R-13/R-14 consideran falta de acompañamiento en pantalanes/varadero y fines de semana. SD7 prevé cuatro semanas de estabilización E1 y estabilización E2 dentro de Operación; la carga parcial de datos no acredita cuadrantes, relevos ni presencia por sitio/turno. Implantación debe conciliar dotación inicial/decreciente por ola, capacidad de fin de semana y atención reforzada con T-15, conservando esta brecha. La formación estacional se prepara en noviembre, antes de habilitar cada persona y fuera del congelamiento, con evaluación de tarea y recuperación. Renovación de más de la mitad de estacionales y todos los monitores, o rechazo por guantes, lluvia y fricción, exige repetir formación y prueba de uso según la sección~\ref{sec:sd8-operacionales}; no basta asistencia ni una interfaz demostrada en oficina. \parencites[RT-20.05--06; p.~35]{sd8-bt}[RT-22.04; numerales 13.3 y 17.6]{sd8-c07}
